\documentclass[10pt,a4paper]{amsart}
\usepackage[margin=19mm]{geometry}
\usepackage{amsmath,amssymb,mathtools}
\usepackage[scr=rsfs,cal=euler]{mathalfa}
\usepackage{xfrac}
\usepackage[unicode,hidelinks,bookmarks=false]{hyperref}
\newcommand{\ie}{i.e.\ }
\newcommand{\cf}{cf.\ }
\newcommand{\resp}{respectively\ }
\newcommand{\Subsection}{\subsection}
\providecommand{\nobreakdash}{\nobreak\hbox{-}}
\setlength{\emergencystretch}{2em}
\multlinegap0pt
\usepackage{xcolor}
\usepackage{tikz}
\usepackage{bm}
\usepackage{algorithm}\usepackage{algpseudocode}
\usepackage{amssymb}
\newtheorem{lem}{Lemma}
\def\RR{\mathbb{R}}
\def\aalpha{\boldsymbol{\alpha}}
\def\bbeta{\boldsymbol{\beta}}
\def\cA{\mathcal{A}}
\def\ee{\mathbf{e}}
\def\ff{\mathbf{f}}
\title{A reduced-order model for parametrized Optimal Transport problems}
\author{Elise Bonnet-Weill, Virginie Ehrlacher, Luca Nenna}
\date{Source: arXiv:2604.09325v1 (CC BY 4.0)}
\begin{document}
\maketitle

\noindent\emph{Source and licence.} A reduced-order model for parametrized Optimal Transport problems. Version arXiv:2604.09325v1 (CC BY 4.0), distributed by arXiv under the Creative Commons Attribution 4.0 International licence, http://creativecommons.org/licenses/by/4.0/, as recorded in the arXiv licence metadata for this exact version and shown on the abstract page https://arxiv.org/abs/2604.09325v1. Full licence notice: \texttt{licenses/arxiv-2604.09325-CC-BY-4.0.txt}.\par\smallskip

\noindent\textit{Editorial scope.} Counted unit: the article's lem points (source0.tex lines 969--971). The article's complete original proof of it is delivered with it (source0.tex lines 973--1033, one \texttt{proof} environment of 61 lines). No mathematics is written, repaired or re-derived: the statement and the proof are the article's own lines, in the article's own notation, and no retained line is substituted or re-flowed.  No further source-local context is needed: every cross-reference of every retained line resolves inside the two delivered spans.  Nothing was omitted from any retained span, so the package carries no omission record.  No retained line contains a citation, so the wrapper carries no bibliography and no bibliography entry.  No retained line contains a figure environment, an image-inclusion command or drawing code, so no image is delivered and the package declares no asset dependency.  Editorial formatting only, in addition: an AMS \texttt{amsart} preamble that restates the article's own declarations and macros used by the retained lines, namely the environments \texttt{lem} on separate counters, together with the standard packages those lines need.  The retained environments print in this document as Lemma 1 in order of appearance, whereas the article prints the same items with the numbering of its own full document.  The complete native source of the article as distributed in the arXiv e-print for this exact version is bundled unchanged as \texttt{sources/198200/source0.tex}.
\begin{lem}\label{lemma: extreme points}
    The set $\displaystyle  \cA := \left\lbrace \aalpha := ({\aalpha^x, \aalpha^y})\in \mathbb{R}_+^{K_x + K_y}, \sum_{k=1}^{K_x} \mathbf{\aalpha}^x_k = 1, \sum_{l=1}^{K_y} \mathbf{\aalpha}^y_l = 1 \right\rbrace$  has  $K_x\times K_y$ extreme points, defined by $\cA_{\text{ext}} = \left\lbrace (\ee_{k}, \ff_{l})\right\rbrace_{\footnotesize \begin{matrix} 1 \leq k \leq K_x \\ 1\leq l \leq K_y \end{matrix}}$, where $\left\lbrace \ee_k \right\rbrace_{1 \leq k \leq K_x}$ and $\left\lbrace \ff_k \right\rbrace_{1 \leq l \leq K_y}$ are the elementary basis vectors of $\RR^{K_x}$ and $\RR^{K_y}$.
\end{lem}

\begin{proof}
We begin by showing the generative property of $\cA_{\text{ext}}$. We then prove that the elements of $\cA_{\text{ext}}$ are extreme points by noticing each point $\aalpha$ of $\cA_{\text{ext}}$, can be expressed as one unique convex combination of points of $\cA_{\text{ext}} $.

Let $(\aalpha^x, \aalpha^y) = \left(\left\lbrace \aalpha_{k}^x \right\rbrace_{1 \leq k \leq K_x}, \left\lbrace \aalpha_{l}^y \right\rbrace_{1\leq l \leq K_y} \right)$ be an element of $\cA$. Let us show that $\aalpha$ is a convex combination of  $\left\lbrace \left( \ee_{k}, \ff_{l}\right) \right\rbrace_{\footnotesize \begin{matrix} 1 \leq k \leq K_x \\ 1\leq l \leq K_y \end{matrix}}$ associated with the coefficients $\bbeta = \left\lbrace \aalpha_{k}^x \aalpha_{l}^y\right\rbrace_{\footnotesize \begin{matrix} 1 \leq k \leq K_x \\ 1\leq l \leq K_y \end{matrix}} $:


\begin{equation}
    \sum_{k' = 1}^{K_x} \sum_{l' = 1}^{K_y} \aalpha_{k'}^x\aalpha_{l'}^y (\ee_{k'}, \ff_{l'})  = (\aalpha^x, \aalpha^y).
\end{equation}

We first check that the coefficients of $\bbeta$ form a convex combination. This arises from the fact that the coefficients of $\aalpha^x$ and $\aalpha^y$ are themselves  convex combinations, so for any $1\leq k' \leq K_x$ and $1 \leq l' \leq K_y$, $\bbeta_{k', l'} \geq 0$, and


\begin{equation}
   \sum_{k' = 1}^{K_x} \sum_{l' = 1}^{K_y} \bbeta_{k', l'} = \sum_{k' = 1}^{K_x} \sum_{l' = 1}^{K_y} \aalpha_{k'}^x\aalpha_{l'}^y  = 1 .
\end{equation}



Now, notice that we can decompose $\aalpha^x$ and $\aalpha^y$ as sums of elementary vectors: $\displaystyle  \aalpha^x = \sum_{k = 1}^{K_x} \aalpha^x_{k} \ee_{k}$, and $\displaystyle  \aalpha^y = \sum_{l = 1}^{K_y} \aalpha^y_{l} \ff_{l}$.

Thus, for $1 \leq k \leq K_x$, 
\begin{equation}
    \begin{array}{ll}
        \displaystyle{} \left(\sum_{k' = 1}^{K_x} \sum_{l' = 1}^{K_y} \aalpha_{k'}^x \aalpha_{l'}^y \ee_{k'}\right)_{k} & = \displaystyle{}
          \left(\sum_{k' = 1}^{K_x} \aalpha_{k'}^x \ee_{k'} \right)_{k}\\
         & \displaystyle{} =  (\aalpha_{k}^x \ee_{k})_{k} \\
         &\displaystyle{} = \aalpha^x_{k},
    \end{array}
\end{equation}

and for $1 \leq l \leq K_y$, 
\begin{equation}
    \begin{array}{ll}
       \displaystyle{}   \left(\sum_{k' = 1}^{K_x} \sum_{l' = 1}^{K_y} \aalpha_{k'}^x \aalpha_{l'}^y  \ff_{l'}\right)_{l} & = 
        \displaystyle{}  \left(\sum_{l' = 1}^{K_y} \aalpha_{l'}^y \ff_{l'} \right)_{l}\\
         &\displaystyle{} = ( \aalpha_{l}^y \ff_{l})_{l} \\
         &\displaystyle{} = \aalpha^y_{l}
    \end{array}
\end{equation}

This concludes the fact that $\left\lbrace \left( \ee_{k}, \ff_{l}\right) \right\rbrace_{\footnotesize \begin{matrix} 1 \leq k \leq K_x \\ 1\leq l \leq K_y \end{matrix}}$ is a generative set of $\cA$.

Now, suppose that, for $1\leq k \leq K_x$ ans $1\leq l \leq K_y$, the point $\displaystyle{} (\ee_{k'}, \ff_{l'})$ is expressed as a convex combination of $\left\lbrace \left( \ee_{k}, \ff_{l}\right) \right\rbrace_{\footnotesize \begin{matrix} 1 \leq k \leq K_x \\ 1\leq l \leq K_y \end{matrix}}$ associated to the coefficients $\{ \bbeta_{k', l'}\}$. Then, in particular :
%
\begin{equation}
    \ee_{k} = \sum_{k' =1}^{K_x} \left( \sum_{l' =1}^{K_y} \bbeta_{k', l'} \right) \ee_{k'},
\end{equation}
%
By linear independence of the elementary vectors $\ee_{k'}$, if $k' \neq k$,
%
 \begin{equation}
     \sum_{l' = 1}^{K_y} \bbeta_{k', l'} = 0,
 \end{equation}
 %
and by non-negativity of the coefficients of $\bbeta$, we obtain that $\bbeta_{k', l' = 0}$ if $k' \neq k$.

By a similar computation, since the $\{\ff_{l'}\}_{1 \leq l'\leq K_y}$ are independent, by non-negativity of the coefficients of $\bbeta$, for every $1\leq k' \leq K_x$, $\leq l' \leq K_y$, $\bbeta_{k', l'} = 0$ if $l' \neq l$. 

Since the coefficients of $\bbeta$ sum to one, this means that the only coefficient that can be non zero, namely $\bbeta_{k,l}$, is equal to 1. So we have proven that every  element of $\left\lbrace \left( \ee_{k}, \ff_{l}\right) \right\rbrace_{\footnotesize \begin{matrix} 1 \leq k \leq K_x \\ 1\leq l \leq K_y \end{matrix}}$ can be written as one unique convex combination of elements of this set.
\end{proof}

\end{document}
